\section{Conclusiones}
\label{sec:conclusiones}

\begin{table}[htbp]
    \centering
    \caption{Conclusiones del an\'{a}lisis de flexibilidad de la l\'{i}nea SIM-009.}
    \label{tab:conclusiones}
    \setcounter{itemcount}{0}
    \begin{tabularx}{\textwidth}{@{}NX@{}}
        \toprule
        \multicolumn{1}{c}{\textbf{Item}} & \textbf{Conclusion} \\
        \midrule
        & La tuber\'{i}a de succi\'{o}n de la bomba de recirculaci\'{o}n 2 del sistema DDW (SIM-009) cumple los criterios de esfuerzo de ASME B31.3-2016 en los nueve casos de carga evaluados: la evaluaci\'{o}n de cumplimiento de c\'{o}digo del modelo arroj\'{o} CODE COMPLIANCE EVALUATION PASSED. \\
        & El esfuerzo de c\'{o}digo m\'{a}ximo es \num{30366.4}~\si{kPa} en el nodo 60, caso 2 (Alt-SUS) W+P2, con una relaci\'{o}n de utilizaci\'{o}n de 26.4~\% frente al admisible de \num{115142.4}~\si{kPa}; la condici\'{o}n que gobierna el dise\~{n}o es la carga sostenida de peso y presi\'{o}n. \\
        & Los desplazamientos m\'{a}ximos del sistema son 10.2~\si{mm} (componente axial $DZ$, caso 3 de operaci\'{o}n con $T_2$) y las cargas en soportes son moderadas (carga vertical m\'{a}xima \num{2382}~\si{N} en el soporte del nodo 60), por lo que la configuraci\'{o}n de soporter\'{i}a del isom\'{e}trico es adecuada. \\
        & La comparaci\'{o}n de las cargas en la boquilla de succi\'{o}n de la bomba con los admisibles del fabricante del equipo queda pendiente como verificaci\'{o}n externa, por no disponerse de dichos admisibles a la fecha. \\
        \bottomrule
    \end{tabularx}
\end{table}

En s\'{i}ntesis, la l\'{i}nea SIM-009 es t\'{e}cnicamente viable seg\'{u}n
ASME B31.3-2016 con un margen de utilizaci\'{o}n cercano al 74~\% del
admisible. No se identifican riesgos de sobre-esfuerzo, interferencias por
expansi\'{o}n ni sobrecarga de soportes, por lo que no se requieren acciones
correctivas sobre la configuraci\'{o}n analizada.
